% Part:sets-functions-relations
% Chapter: size-of-sets
% Section: non-enumerability-alt

\documentclass[../../../include/open-logic-section]{subfiles}

\begin{document}

\olfileid{sfr}{siz}{nen-alt}

\olsection{\printtoken{S}{nonenumerable} ગણો}

\begin{editorial}
  આ વિભાગ \olref[enm-alt]{sec}ની વ્યાખ્યાઓ વાપરીને
  $\Bin^\omega$ અને $\Pow{\Nat}$ની અગણનીયતા સાબિત કરે છે;
  એટલે કે, $\PosInt$થી આવતા વ્યાપ્ત વિધેયને બદલે~$\Nat$ સાથેનું
  એક-એક વ્યાપ્ત વિધેય જરૂરી માને છે.
\end{editorial}

\begin{explain}
પ્રાકૃતિક સંખ્યાઓનો ગણ $\Nat$ અનંત છે. તે તુચ્છ રીતે
!!{enumerable} પણ છે. પરંતુ નોંધપાત્ર હકીકત એ છે કે
\emph{!!{nonenumerable}} ગણો હોય છે, એટલે કે એવા ગણો જે
!!{enumerable} નથી (\olref[sfr][siz][enm-alt]{defn:enumerable} જુઓ).

આ આશ્ચર્યજનક લાગી શકે. ગણ $A$ એ !!{nonenumerable} છે તેમ કહેવાનો
અર્થ આખરે એ છે કે \emph{કોઈ} !!{bijection} $f \colon \Nat \to A$
નથી; એટલે કે, $\Nat$ના અનંત ઘણા !!{element}sનું~$A$માં નિરૂપણ
કરતું કોઈ વિધેય આખા~$A$ને નિઃશેષ આવરી શકતું નથી. આમ, જો $A$ એ
!!{nonenumerable} હોય, તો ગણ~$A$માં પ્રાકૃતિક સંખ્યાઓ કરતાં “વધુ”
!!{element}s હોય છે.

કોઈ ગણ !!{nonenumerable} છે તે સાબિત કરવા, યોગ્ય !!{bijection}
અસ્તિત્વમાં હોઈ ન શકે તે બતાવવું પડે. તેની શ્રેષ્ઠ રીત એ બતાવવાની
છે કે $A$ના !!{element}sની પરિગણનાનો દરેક પ્રયત્ન ઓછામાં ઓછો એક
!!{element} છોડી જ દે છે; તેનાથી કોઈ વિધેય $f\colon \Nat \to A$
!!{surjective} નથી તેમ મળે છે. આ સ્થાપિત કરવાની એક સામાન્ય યુક્તિ
કૅન્ટૉરની \emph{વિકર્ણ પદ્ધતિ} વાપરવાની છે. $A$ના !!{element}sની
કોઈ યાદી, જેમ કે $x_1$, $x_2$, \dots, આપેલી હોય, તો આપણે $A$નો
એવો બીજો !!{element} બનાવીએ છીએ જે તેની રચનાને કારણે આ યાદીમાં
હોઈ જ ન શકે.

આ બધું ઉદાહરણથી સૌથી સારી રીતે સમજાય છે. તેથી આપણો પહેલો દાખલો
$0$ અને $1$ની બધી અનંત પ્રતીકશ્રેણીઓનો ગણ~$\Bin^\omega$ છે.
(“$\Bin$” દ્વિઆધારી માટે છે, અને આપણે તેને માત્ર બે ઘટકોવાળો ગણ
$\{0,1\}$ માની શકીએ.)\oliflabeldef{sfr:card-arithmetic:card-opps:sec}{\footnote{વધુ
ચોક્કસ રીતે, $\Bin^\omega$ એ $0$ અને $1$ના બધા $\omega$-અનુક્રમોનો
ગણ, એટલે કે $\funfromto{\omega}{\{0,1\}}$, છે એવું ઠરાવવું જોઈએ.
પરંતુ તેનો અર્થ \olref[sfr][card-arithmetic][card-opps]{sec}માં જ
સ્પષ્ટ થશે.}}{}
જોકે હાલ આ સહેજ ઢીલી રજૂઆતથી કોઈ ગૂંચવણ થવી જોઈએ નહીં.
\sourcecorrection{OLSIZ-011}{મૂળની \texttt{non-enumerability-alt.tex},
પંક્તિઓ 47--53માં \texttt{oliflabeldef}નો ત્રીજો અવયવ
બંધ થતો નથી: ખાલી વૈકલ્પિક શાખા અને પછીનું સામાન્ય વાક્ય બંને
બીજી શાખાની અંદર રહી જાય છે. અહીં પાદનોંધને હકારાત્મક શાખા,
ખાલી ભાગને નકારાત્મક શાખા અને પછીના વાક્યને સશરતી આદેશની બહાર રાખ્યાં છે.}
\end{explain}

\begin{thm}
\ollabel{thm:nonenum-bin-omega}
$\Bin^\omega$~એ !!{nonenumerable} છે.
\end{thm}

\begin{proof}
$\Bin^\omega$ના કોઈ ઉપગણની મનસ્વી પરિગણના વિચારો. એટલે આપણી પાસે
કોઈ યાદી $s_{0}$, $s_{1}$, $s_{2}$, \dots{} છે, જેમાં દરેક $s_n$
$0$ અને~$1$ની અનંત પ્રતીકશ્રેણી છે. $s_n(m)$ને આ યાદીની $n$મી
પ્રતીકશ્રેણીનો $m$મો અંક કહીએ. હવે આપણી યાદીને એવી સરણિ તરીકે
વિચારી શકાય જેમાં $s_n(m)$ને $n$મી હાર અને $m$મા સ્તંભમાં મૂક્યું છે:
\sourcecorrection{OLSIZ-008}{મૂળની \texttt{non-enumerability-alt.tex},
પંક્તિઓ 59--62માં \(s_n(m)\)ને “\(m\)મી પ્રતીકશ્રેણીનો \(n\)મો અંક”
કહ્યો છે, જ્યારે તરત પછીનું સ્થાન અને કોષ્ટક તેને \(n\)મી
પ્રતીકશ્રેણીનો \(m\)મો અંક બતાવે છે. અહીં કોષ્ટક મુજબ લખ્યું છે.}
\[
\begin{array}{c|c|c|c|c|c}
& 0 & 1 & 2 & 3 & \dots \\\hline
0 & \mathbf{s_{0}(0)} & s_{0}(1) & s_{0}(2) & s_0(3) & \dots \\\hline
1 & s_{1}(0)& \mathbf{s_{1}(1)} & s_1(2) & s_1(3) & \dots \\\hline
2 & s_{2}(0)& s_{2}(1) & \mathbf{s_2(2)} & s_2(3) & \dots \\\hline
3 & s_{3}(0)& s_{3}(1) & s_3(2) & \mathbf{s_3(3)} & \dots \\\hline
\vdots & \vdots & \vdots & \vdots & \vdots & \mathbf{\ddots}
\end{array}
\]
હવે આપણે $0$ અને $1$ની એવી અનંત પ્રતીકશ્રેણી~$d$ બનાવીશું જે આ
યાદીમાં નથી. તેના દરેક પદને નિર્દિષ્ટ કરીને, એટલે કે બધા
$n \in \Nat$ માટે $d(n)$ નિર્દિષ્ટ કરીને, તેમ કરીશું. સહજ રીતે,
ઉપરની સરણિના વિકર્ણને નીચે તરફ વાંચીને (તેથી “વિકર્ણ પદ્ધતિ” નામ)
દરેક $1$ને $0$માં અને દરેક $0$ને~$1$માં બદલીને તેમ કરીએ છીએ.
\sourcecorrection{OLSIZ-009}{મૂળની \texttt{non-enumerability-alt.tex},
પંક્તિઓ 76--78માં બંને વાર “દરેક \(1\)ને \(0\)માં” લખ્યું છે.
તરત પછીની કિસ્સાવાર વ્યાખ્યા મુજબ બીજી ક્રિયા દરેક \(0\)ને \(1\)માં
બદલવાની છે; અહીં તે મુજબ લખ્યું છે.}
વધુ અમૂર્ત રીતે, વિકર્ણનો $n$મો !!{element}~$s_n(n)$ એ $1$ છે કે
$0$ તે પ્રમાણે $d(n)$ને $0$ કે $1$ વ્યાખ્યાયિત કરીએ છીએ; એટલે કે:
\[
d(n) =
\begin{cases}
1 & \text{જો $s_{n}(n) = 0$ હોય}\\
0 & \text{જો $s_{n}(n) = 1$ હોય}
\end{cases}
\]
આ $0$ અને $1$ની અનંત પ્રતીકશ્રેણી હોવાથી સ્પષ્ટ છે કે
$d \in \Bin^\omega$. પરંતુ આપણે $d$ને એવી રીતે બનાવ્યો છે કે
કોઈ પણ $n \in \Nat$ માટે $d(n) \neq s_n(n)$. એટલે કે, $d$ તેના
$n$મા પદે $s_n$થી જુદો છે. તેથી કોઈ પણ $n\in \Nat$ માટે
$d \neq s_n$. આમ $d$ યાદી $s_0$, $s_1$, $s_2$, \dots{}માં હોઈ ન શકે.

$\Bin^\omega$ના કોઈ ઉપગણની મનસ્વી પરિગણના આપેલી હોય, તો તેમાં
$\Bin^\omega$નો કોઈ !!{element} રહી જાય છે તે આપણે બતાવ્યું છે.
તેથી ગણ $\Bin^\omega$ની કોઈ પરિગણના નથી; એટલે $\Bin^\omega$ એ
!!{nonenumerable} છે.
\end{proof}

\begin{explain}
આ સાબિતીની રીતને “વિકર્ણીકરણ” કહે છે, કારણ કે તે $d$ને વ્યાખ્યાયિત
કરવા સરણિના વિકર્ણનો ઉપયોગ કરે છે. જોકે, વિકર્ણીકરણમાં સરણિ હાજર
હોય જ એવું જરૂરી નથી. ખરેખર, કોઈ સરણિ અને કોઈ વાસ્તવિક વિકર્ણ ન
હોય ત્યારે પણ સમાન વિચારથી કોઈ ગણ !!{nonenumerable} છે તે બતાવી
શકાય છે. નીચેનું પરિણામ તે કેવી રીતે થાય છે તે દર્શાવે છે.
\end{explain}

\begin{thm}
\ollabel{thm:nonenum-pownat}
$\Pow{\Nat}$ એ !!{enumerable} નથી.
\end{thm}

\begin{proof}
એ જ રીતે આગળ વધીએ અને બતાવીએ કે $\Nat$ના ઉપગણોની દરેક યાદીમાં
$\Nat$નો કોઈ ઉપગણ રહી જાય છે. ધારો કે $\Nat$ના ઉપગણોની કોઈ યાદી
$N_0, N_1, N_2, \ldots$ છે. ગણ $D$ને આ રીતે વ્યાખ્યાયિત કરીએ:
$n \in D$ જો અને માત્ર જો $n \notin N_{n}$:
\[
D = \Setabs{n \in \Nat}{n \notin N_n}
\]
સ્પષ્ટ છે કે $D\subseteq \Nat$. પરંતુ $D$ યાદીમાં હોઈ ન શકે.
કારણ કે રચના પ્રમાણે $n \in D$ જો અને માત્ર જો $n\notin N_n$,
તેથી કોઈ પણ $n \in \Nat$ માટે $D \neq N_n$.
\end{proof}

\begin{explain}
અગાઉની સાબિતીમાં વિકર્ણનો ઉલ્લેખ થયો ન હતો. છતાં તેને આ રીતે
ચિત્રિત કરીને તેમાં વિકર્ણ હોય તેમ વિચારી શકો: ગણો $N_0$, $N_1$,
\dots{}ને સરણિમાં લખેલા કલ્પો, જ્યાં $N_n$ને $n$મી હારમાં એ રીતે
લખીએ કે $m \in N_n$ હોય જો અને માત્ર જો $m$ને $m$મા સ્તંભમાં
લખીએ. દાખલા તરીકે, ધારો કે યાદીના પહેલા ચાર ગણો
$\{0,1,2,\dots\}$, $\{1, 3, 5, \dots\}$, $\{0,1,4\}$ અને
$\{2,3,4,\dots\}$ છે; ત્યારે આપણી સરણિની શરૂઆત આવી હશે:
\[
\begin{array}{r@{}rrrrrrr}
  N_0 = \{ & \mathbf{0}, & 1, & 2, & & & & \dots\}\\
  N_1 = \{ &  & \mathbf{1}, &  & 3, &  & 5, & \dots\}\\
  N_2 = \{ & 0, & 1, &  &  & 4\phantom{,} &  & \}\\
  N_3 = \{ &  &  & 2, & \mathbf{3}, & 4, & & \dots\}\\
  &\vdots & & & & & & \ddots\phantom{\}}
  \end{array}
\]
પછી વિકર્ણ પર નીચે જતાં, $n$ વિકર્ણ પર \emph{નથી} જો અને માત્ર જો
$n \in D$ રાખીને મળતો ગણ એ $D$ છે. તેથી ઉપરના કિસ્સામાં આપણે
$0$ અને $1$ને છોડી દઈશું,~$2$નો સમાવેશ કરીશું,~$3$ને છોડીશું, વગેરે.
\end{explain}

\begin{prob}
સ્પષ્ટ વિકર્ણ દલીલ વડે બતાવો કે બધા વિધેયો
$f \colon \Nat \to \Nat$નો ગણ !!{nonenumerable} છે. એટલે કે,
જો $f_1$, $f_2$, \dots, વિધેયોની યાદી હોય અને દરેક
$f_i\colon \Nat \to \Nat$ હોય, તો આ યાદીમાં ન હોય એવું કોઈ
$g \colon \Nat \to \Nat$ છે તે બતાવો.
\end{prob}

\end{document}
